\clearpage
\fancyhead[L]{\small CONSTANTS OF THE BIFILAR MAGNETOMETER.}
\SourceLeaf{410}\pdfbookmark[0]{On determining the constants of the bifilar magnetometer.}{bifilarconstanten-title}
\begin{center}
{\Large\bfseries ON DETERMINING\par\medskip THE CONSTANTS OF THE BIFILAR MAGNETOMETER.}
\par\bigskip
Resultate aus den Beobachtungen des magnetischen Vereins. 1840. I.
\end{center}
\renewcommand{\Article}[1]{\par\medskip\Needspace{4\baselineskip}\pdfbookmark[1]{#1.}{bifilarconstanten-art-#1}\centerline{\textbf{#1.}}\nobreak\smallskip}%
\Article{1}
To use the bifilar magnetometer correctly and reliably, we need to know the numerical values of certain quantities. They relate to arrangements of the parts of the instrument that may be regarded as constant under certain conditions. These quantities are essential to determining the equilibrium positions and swing times observed at different settings of the movable parts. There are four such quantities.

1) The position that the index of the mirror alidade must have for the normal to the mirror and the optical axis of the observing telescope to lie in one vertical plane when the two suspension wires are in a vertical plane. This position, with the reflecting face of the mirror turned towards the telescope, is to be denoted by $P$.

2) The position that must be given to the index of the cradle, with the suspension wires in the same position, for the magnetic axis of the bar magnet to be in its natural position in the magnetic meridian. I denote this position by $Q$.

It goes without saying that if each alidade has more than one index, one of them must always be chosen at will as the principal index.

3) The ratio of the magnetic directive force to the one arising from the manner of suspension, which we may call the static directive force. This ratio is to be expressed by $R:1$.

\SourceLeaf{411}4) The static swing duration of the apparatus, that is, the duration that would result from the manner of suspension alone, with no action of the Earth’s magnetism on the bar magnet. I denote the square of this swing duration by $S$.

It follows that $\frac{S}{R}$ expresses the square of the purely magnetic swing duration, that is, the duration that would occur if the apparatus were suspended from a single thread without torsion.

\Article{2}
We must first work out how the direct readings of the bifilar magnetometer are related to the positions of the two alidades and to these four quantities.

With the mirror alidade at $A$ and the cradle alidade at $B$, let $t$ denote the swing duration and $p$ the distance, converted into divisions of arc, from the scale mark corresponding to the equilibrium position to the point on the scale that lies in the same vertical plane as the optical axis of the observing telescope. This point is marked by the plumb line hanging from the centre of the objective. To make the directions clear, I assume that the divisions on the circle and on the scale run from left to right, and take $p$ as positive when the number seen at the telescope’s crosshairs is greater than the number at the plumb line. In that equilibrium position, then, the bifilar magnetometer is turned back through $A-P-p$, that is, from right to left relative to the position in which the suspension wires were parallel. Equivalently, $A-P-p$ is the angle between the straight line through the two lower ends of the suspension wires and a line parallel to the one joining their two upper ends. The torque produced by the manner of suspension is only approximately proportional to the sine of this angle, but the approximation is accurate enough for practical use. We put it $=D\sin(A-P-p)$, where $D$ denotes the static directive force. Positive values of the torque refer to rotation from left to right.

In the same position of the apparatus the magnetic axis of the bar magnet makes with the magnetic meridian the angle cou%
\expandafter\def\csname HistoricalNoteText416\endcsname{The magnetism of the needle of a compass on a sundial made in 1709 and kept at the observatory here could in 1841 be increased only threefold by stroking the needle anew to saturation. In another from 1603 it could be increased to only five times its former strength. On the very probable assumption that neither had been stroked again since it was made, and if we also assume that both had originally been magnetized to saturation and that their strength had gradually decreased in geometric progression, the annual loss amounts to $\frac{1}{120}$ in the former and $\frac{1}{150}$ in the latter. It is even smaller if the original magnetization fell short of saturation.}%
\SourceLeaf{412}nted from right to left, $A-P-p-B+Q$. The torque due to the action of the Earth’s magnetism on the bar magnet, reckoned positive from left to right, is $=RD\sin(A-P-p-B+Q)$. Hence we have equation (1)
\[
0=\sin(A-P-p)+R\sin(A-P-p-B+Q)
\]
If the whole apparatus is turned from its equilibrium position through the angle $z$ from right to left, the torque acting in the opposite direction is
\[
D\sin(z+A-P-p)+DR\sin(z+A-P-p-B+Q)
\]
Expanding the two sines and taking equation (1) into account gives
\[
D\sin z(\cos(A-P-p)+R\cos(A-P-p-B+Q))
\]
which is proportional to the sine of $z$. We must therefore regard
\[
D(\cos(A-P-p)+R\cos(A-P-p-B+Q))
\]
as the directive force resulting from the combination of the static and magnetic directive forces, so we have (2)
\[
\frac{S}{tt}=\cos(A-P-p)+R\cos(A-P-p-B+Q)
\]
Squaring both sides of each of equations (1) and (2) and adding gives (3)
\[
\frac{SS}{t^4}=1+2R\cos(Q-B)+RR
\]
If $e$ denotes the base of the hyperbolic logarithms and $i$ the imaginary unit $\BareRoot{-1}$, the two equations (1) and (2) can conveniently be combined into one
\[
\frac{S}{tt}=e^{i(A-P-p)}+Re^{i(A-P-p-B+Q)}
\]
or, more simply still, into the following equation (4)
\[
1=\frac{S}{tt}e^{i(P+p-A)}-Re^{i(Q-B)}
\]
\SourceLeaf{413}which comprises the two equations
\[
\begin{gathered}
1=\frac{S}{tt}\cos(P+p-A)-R\cos(Q-B)\\
\frac{S}{tt}\sin(P+p-A)=R\sin(Q-B)
\end{gathered}
\]

For the natural position, where $Q=B$, we obtain
\[
tt=\frac{S}{1+R}
\]
and for the reversed position, where $Q=B+180^0$,
\[
tt=\frac{S}{1-R}
\]

\Article{3}
The transverse position, in the narrower sense, requires
\[
A-P-p-B+Q=\pm90^0
\]
where the upper sign applies when the north pole of the bar magnet is to lie to the west of the magnetic meridian, and the lower to the eastward position. From (1) we thus obtain
\[
\sin(A-P-p)=\mp R
\]
If $\varphi$ denotes the acute angle whose sine is $=R$, then for the westward position of the north pole we have
\[
A=P+p-\varphi,\quad B=Q-\varphi-90^0
\]
and for the eastward position
\[
A=P+p+\varphi,\quad B=Q+\varphi+90^0
\]
For the equilibrium position to correspond to the scale mark indicated by the plumb line itself, the mirror alidade must therefore be set to $P-\varphi$ in the first case and to $P+\varphi$ in the second.

For the swing duration corresponding to the transverse position, formula (2) immediately gives
\SourceLeaf{414}\[
\frac{S}{tt}=\cos\varphi
\]
or
\[
tt=\frac{S}{\BareRoot{(1-RR)}}
\]
The swing duration in the transverse position is therefore the geometric mean of the swing times in the natural and reversed positions.

\Article{4}
To see clearly to what extent the quantities can vary, we must trace them to their origins.

The angles $P$ and $Q$ each consist of three parts. The first part of $P$ is the angle between the radius to the zero point of the circle and the straight line joining the two lower ends of the suspension wires. The second is the angle between the straight line joining the two upper ends of the suspension wires and the optical axis of the telescope (or rather between the projections of this straight line onto a horizontal plane, which is also tacitly understood for all other sides of angles that are not themselves horizontal or do not intersect one another directly). The third is the angle between the normal to the mirror and the radius to the principal index of the mirror alidade.

The first part of $Q$ is the same as the first part of $P$. The second is the angle between the straight line joining the two upper ends of the suspension wires and the magnetic meridian. The third is the angle between the magnetic axis of the bar magnet in the cradle and the radius to the principal index of the cradle alidade.

All five angles are to be counted from left to right. From this analysis it is clear that as long as the suspension of the instrument, the attachment of the mirror to its alidade and the position of the observing telescope remain fixed, $P$ will be entirely invariable. Because of its second part, however, $Q$ undergoes precisely the same changes as the magnetic declination, changes from left to right being counted as positive.

\SourceLeaf{415}The static directive force is expressed by the formula
\[
D=\frac{fgG}{4h}
\]
where $G$ denotes the weight of the apparatus (that is, the mass multiplied by gravity), $f$ the distance between the suspension wires at their lower ends, $g$ that distance at their upper ends, and $h$ the height of the upper attachment above the lower. This holds at least if we set aside the small increase that the force gains from the resistance of the individual suspension wires to torsion. We can readily do so here, where our immediate concern is the variability of the total force. If we also denote by $K$ the moment of inertia about the vertical axis of rotation, then, with $\pi$ in its usual sense,
\[
S=\frac{\pi\pi K}{D}=\frac{4\pi\pi hK}{fgG}
\]
It is clear that the individual factors $f,g,h,K$ vary with changes in temperature. These variations are in part very small in themselves and, as will be shown below, almost completely compensate for one another in the value of $S$. The difference arising from the unequal loss of weight due to unequal air density may be regarded as quite imperceptible.

The magnetic directive force is $=TM$ if $T$ denotes the intensity of the horizontal component of the Earth’s magnetism and $M$ the magnetic moment of the bar magnet. We therefore have
\[
R=\frac{TM}{D}=\frac{4TMh}{fgG}=\frac{STM}{\pi\pi K}
\]

The variability of $R$ thus has three causes.

The first is the continual variability of $T$. The second is variation in temperature, which affects the linear quantities $f,g,h$ as well as the magnetism $M$ of the bar. The third is the variability of $M$ independently of the temperature at any given time.

Our knowledge of the third cause is still rather incomplete. In the experiments reported in volume 2 of the \textit{Resultate} by Professor \textsc{Weber}, the loss due to artificial heating was never fully recovered by subsequent cooling.
\SourceLeaf{416}A substantial lasting loss remained after the initial temperature was restored. On the other hand, experience shows that magnetic needles retain a considerable degree of magnetism for many years without further stroking, despite daily and annual changes in temperature. From this we must infer an extremely slow progressive loss\HistoricalNote{416}. It would be very important to know precisely the conditions under which a change in temperature produces the smallest possible lasting loss of magnetic strength. In addition to the nature and hardening of the steel and strong initial magnetization, the main requirements are likely to be that some time has elapsed since that magnetization, that the temperature changes stay within certain limits, and that they always occur very slowly and gradually. Under those conditions it should be permissible to treat the magnetic state of a bar magnet (meaning its magnetic moment reduced to a specified reference temperature) as constant for a moderate time, say a few days. If a definite decrease is found after a longer period, we may assume for the interval a steady decrease in geometric progression. Putting into practice Professor \textsc{Weber}’s ingenious proposal in the paper that follows below seems particularly well suited to shedding light on this interesting subject.

\Article{5}
To give a precise meaning to the task of determining the numerical values of the quantities of a bifilar magnetometer by experiment, we take the values sought for the variable quantities to be those corresponding to a specified declination, a specified horizontal intensity, a specified temperature, and the magnetic state of the
\SourceLeaf{417}bar magnet at the time of the experiments. Changes that the bar may undergo after a longer interval are therefore outside the question. We denote these reference values of the variable quantities by $Q^0,R^0,S^0$ (since $P$ is already constant in itself), and in general set
\[
Q=Q^0+q,\quad R=rR^0,\quad S=sS^0
\]
Similarly, let $f^0,g^0,h^0,K^0,T^0,M^0$ denote the reference values of the variable quantities $f,g,h,K,T,M$. We thus have at once
\[
s=\frac{f^0g^0hK}{fgh^0K^0},\qquad r=\frac{f^0g^0hMT}{fgh^0M^0T^0}
\]

To allow for changes in the direction and strength of the Earth’s magnetic force during the operations needed to determine these quantities, an auxiliary instrument must of course be available. A unifilar magnetometer is best, with the swing duration and position observed at the same time. This auxiliary magnetometer also makes it possible to identify the reference declination and reference intensity to be chosen. First we make the former correspond to a specified mark on the scale and the latter to a specified swing duration at the reference temperature. This also leaves us free to express both reference quantities in absolute measure by known methods. It follows directly that $q$ is the difference between the position observed on the auxiliary magnetometer and the reference position, converted into divisions of arc. Further, if the quantities denoted by $M,K,t$ for the bifilar magnetometer are denoted by $m,k,\theta$ for the auxiliary magnetometer, with reference values $m^0,k^0,\theta^0$, then
\[
\frac{\theta\theta}{\theta^0\theta^0}=\frac{km^0T^0}{k^0mT}
\]
and consequently
\[
r=\frac{f^0g^0hkm^0M\theta^0\theta^0}{fgh^0k^0mM^0\theta\theta}=\frac{skK^0m^0M\theta^0\theta^0}{k^0KmM^0\theta\theta}
\]

Of the seven factors $\frac{f^0}{f},\frac{g^0}{g},\frac{h}{h^0},\frac{k}{k^0},\frac{K^0}{K},\frac{m^0}{m},\frac{M}{M^0}$ occurring in the expressions for $s$ and $r$, the first five must be calculated from the expansion of the metals concerned under changes in temperature, and the last two from the best knowledge available of the effect of tem\SourceLeaf{418}perature on the magnetism of the bars. What we have called the magnetic state is treated as constant in both bar magnets during the operations at issue here. We add some further details concerning this calculation.

Choosing the freezing point as the reference temperature, we denote by $c$ and $c'$ the temperatures inside the cases of the bifilar and auxiliary magnetometers, and by $c''$ the temperature at the upper attachment of the suspension wires of the former. For a rise of one degree in temperature we further denote the expansion of steel by $\alpha$, that of brass by $\beta$, and the decrease in the magnetism of the bars in the two instruments by $\gamma$ and $\gamma'$. Since by far the greater part of the change in the moment of inertia of each instrument is due to the expansion of the bar magnets themselves, we can set without hesitation
\[
\frac{K}{K^0}=(1+\alpha c)^2,\qquad \frac{k}{k^0}=(1+\alpha c')^2
\]
For the expansion of the suspension wires, if they are steel wires as in the apparatus here, we may retain the same coefficient $\alpha$ and take their temperature to be $\tfrac{1}{2}(c+c'')$, so that
\[
\frac{h}{h^0}=1+\tfrac{1}{2}\alpha(c+c'')
\]
We therefore have (1)
\[
s=\frac{(1+\alpha c)^2(1+\tfrac{1}{2}\alpha(c+c''))}{(1+\beta c)(1+\beta c'')}
\]
which may also be written with sufficient accuracy as
\[
s=1+(3\alpha-2\beta)c-(\beta-\tfrac{1}{2}\alpha)(c''-c)
\]
Since experience shows that the relation $\beta=\tfrac{3}{2}\alpha$ holds to a close approximation, we have, to a close approximation, (2)
\[
s=1-\alpha(c''-c)
\]
That is, the change in $S$ depends only on the difference between the lower and upper temperatures, so that as a rule $S$ may be regarded as entirely constant.

We have further (3)
\[
r=\frac{s\theta^0\theta^0}{\theta\theta}.\frac{1-\gamma c}{1-\gamma'c'}.\left(\frac{1+\alpha c'}{1+\alpha c}\right)^2
\]
\SourceLeaf{419}or, if the temperature changes have the same effect on both bars, that is, if $\gamma'=\gamma$, with sufficient accuracy,
\[
r=\frac{s\theta^0\theta^0}{\theta\theta}(1+(\gamma+2\alpha)(c'-c))
\]
or, in accordance with (2), with the same accuracy (4)
\[
r=\frac{\theta^0\theta^0}{\theta\theta}(1+(\gamma+2\alpha)(c'-c)-\alpha(c''-c))
\]

Finally, it should be noted that the comparison observations made with the unifilar magnetometer give only the mean value of $\frac{T}{T^0}$ over the whole time covered by the swing observations, rather than the value at any particular instant. It goes without saying, then, that all the other quantities with which those swing observations are to be combined, directly or indirectly, as simultaneous must likewise be mean values over the same interval.

\Article{6}
The simplest way to determine the four quantities is as follows.

With the cradle in any position, place a nonmagnetic bar of about the same weight in it in place of the bar magnet. Set the mirror so that, at equilibrium, the image of some point on the scale appears at the crosshairs of the observing telescope. Then, with $A$ and $p$ in the senses given above, $P=A-p$. To make the result less dependent on a very precise knowledge of the value of a scale division or on a very accurate conversion into divisions of arc, one may repeat the operation with a new mirror setting very close to the desired one, if $p$ was still very large the first time. A wooden bar weighted with lead is best suited for this operation. Approximately the same weight is required so that any slight torsion that the suspension wires might have on their own when the whole apparatus is in equilibrium has no effect.

Without changing the mirror setting any further, put the bar magnet in the cradle in place of the previous load. The cradle is then to be set
\SourceLeaf{420}so that the position of rest corresponds to the same point on the scale as it last did with the nonmagnetic load. To do this, find by trial two different cradle settings between which the required setting lies, and apply a simple interpolation procedure to the scale readings at those settings. The bar magnet may be placed in either its natural or its reversed position. In the first case the resulting value of $B$ (the setting of the cradle alidade) is $=Q$, and in the second $=Q\pm180^0$. The reversed position has the advantage of greater precision, since a small change in $B$ corresponds to a large change in the scale reading. The natural position is somewhat more convenient because it is easier to set the cradle without going beyond the limits of the scale. To avoid laborious trial and error, it is therefore best to start in the natural position and use the result only as a preparation for making settings at two closely spaced scale marks in the trials with the bar reversed.

The result for $Q$ refers to the position of the magnetic meridian during or between the last two trials. Without an auxiliary instrument for comparison observations, we can expect only a determination subject to this variation. If an auxiliary instrument is available, however, simultaneous readings of its position give the values of $q$ corresponding to those two observations. Applying the interpolation procedure above to the two values of $B-q$ then gives $Q^0$ or $Q^0\pm180^0$.

Finally, observe the swing duration in both the natural and the reversed positions. For this purpose, set the cradle alidade as accurately as possible to the value of $Q$ (or of $Q+180^0$ in the corresponding case) that holds at the start of the swing observations. Let the swing duration be $t'$ in the natural position and $t''$ in the reversed one. If swings can be observed at the same time on the auxiliary magnetometer, they give the corresponding values of $r$, which may be denoted by $r',r''$. If the variability of $S$ is also to be taken into account, let $s',s''$ be the corresponding values of $s$. The small changes in the position of the magnetic meridian during the swing observations
\SourceLeaf{421}will as a rule have no appreciable effect on the results. The two equations at the end of Article 2 thus become
\[
\begin{aligned}
t't'&=\frac{s'S^0}{1+r'R^0}\\
t''t''&=\frac{s''S^0}{1-r''R^0}
\end{aligned}
\]
from which elimination gives
\[
\begin{aligned}
R^0&=\frac{s't''t''-s''t't'}{r''s't''t''+r's''t't'}\\
S^0&=\frac{(r'+r'')t't't''t''}{r''s't''t''+r's''t't'}
\end{aligned}
\]

Following the remark in Article 5, we may readily treat $S$ as constant, or set $s'=s''=1$, in which case the formulas become
\[
\begin{aligned}
R^0&=\frac{t''t''-t't'}{r''t''t''+r't't'}\\
S&=\frac{(r'+r'')t't't''t''}{r''t''t''+r't't'}
\end{aligned}
\]
If comparison observations with an auxiliary instrument are unavailable, we have no choice but to set directly
\[
\begin{aligned}
R&=\frac{t''t''-t't'}{t''t''+t't'}\\
S&=\frac{2t't't''t''}{t''t''+t't'}
\end{aligned}
\]
It is clear that the value of $R$ thus obtained is only a kind of mean of the values that hold for the two sets of swing observations, while $S$ will retain a slight error that depends on the inequality of those sets.

\Article{7}
We base the more general solution of our problem on simultaneous observations of the swing duration and equilibrium position of the bifilar magnetometer at any two different cradle settings. We denote the particular values of the quantities $A,B,p,Q,R,S,t$

for the first set of observations by $A',B',p',Q^0+q',r'R^0,s'S^0,t'$,

for the second set by $A'',B'',p'',Q^0+q'',r''R^0,s''S^0,t''$.

\SourceLeaf{422}Substituting these values in the two equations (1) and (2) of Article 2 gives four equations from which the unknown quantities $P,Q^0,R^0,S^0$ could be found by elimination. We reach the same goal much more easily by calculating with imaginary quantities. Using formula (4) of Article 2, we start from the two equations
\[
\begin{aligned}
1&=\frac{s'S^0}{t't'}e^{i(P+p'-A')}-r'R^0e^{i(Q^0+q'-B')}\\
1&=\frac{s''S^0}{t''t''}e^{i(P+p''-A'')}-r''R^0e^{i(Q^0+q''-B'')}
\end{aligned}
\]
which, if for brevity we set
\[
\begin{aligned}
\frac{s'}{t't'}e^{i(p'-A')}&=a'\\
\frac{s''}{t''t''}e^{i(p''-A'')}&=a''\\
r'e^{i(q'-B')}&=b'\\
r''e^{i(q''-B'')}&=b''\\
S^0e^{iP}&=x\\
R^0e^{iQ^0}&=y
\end{aligned}
\]
become
\[
\begin{aligned}
1&=a'x-b'y\\
1&=a''x-b''y
\end{aligned}
\]
from which we obtain
\[
\begin{aligned}
x&=\frac{b''-b'}{a'b''-a''b'}=\frac{\dfrac{b''}{b'}-1}{\dfrac{b''}{b'}-\dfrac{a''}{a'}}.\frac{1}{a'}\\[1ex]
y&=\frac{a''-a'}{a'b''-a''b'}=\frac{\dfrac{a''}{a'}-1}{\dfrac{b''}{b'}-\dfrac{a''}{a'}}.\frac{1}{b'}
\end{aligned}
\]
This gives the following detailed rules for the calculation. Set (I)
\SourceLeaf{423}
\[
\begin{aligned}
\frac{t't'}{s'}.\frac{s''}{t''t''}.\cos(A'-A''-p'+p'')&=\mathfrak{A}\\
\frac{t't'}{s'}.\frac{s''}{t''t''}.\sin(A'-A''-p'+p'')&=\mathfrak{A}_1\\
\frac{r''}{r'}.\cos(B'-B''-q'+q'')&=\mathfrak{B}\\
\frac{r''}{r'}.\sin(B'-B''-q'+q'')&=\mathfrak{B}_1
\end{aligned}
\]
so that
\[
\begin{aligned}
\frac{a''}{a'}&=\mathfrak{A}+i\mathfrak{A}_1\\
\frac{b''}{b'}&=\mathfrak{B}+i\mathfrak{B}_1
\end{aligned}
\]
Then determine the six quantities $u,U,v,V,w,W$ from the equations (II)
\[
\begin{aligned}
\mathfrak{A}-1&=u\cos U\\
\mathfrak{A}_1&=u\sin U\\
\mathfrak{B}-1&=v\cos V\\
\mathfrak{B}_1&=v\sin V\\
\mathfrak{B}-\mathfrak{A}&=w\cos W\\
\mathfrak{B}_1-\mathfrak{A}_1&=w\sin W
\end{aligned}
\]
with $u,v,w$ chosen to be positive. Then
\[
\begin{aligned}
\frac{b''}{b'}-1&=ve^{iV}\\
\frac{a''}{a'}-1&=ue^{iU}\\
\frac{b''}{b'}-\frac{a''}{a'}&=we^{iW}
\end{aligned}
\]
and consequently
\[
\begin{aligned}
x&=\frac{t't'}{s'}.\frac{v}{w}e^{i(V-W+A'-p')}\\
y&=\frac{u}{r'w}e^{i(U-W+B'-q')}
\end{aligned}
\]
from which one easily concludes that (III)
\SourceLeaf{424}
\[
\begin{aligned}
P&=V-W+A'-p'\\
Q^0&=U-W+B'-q'\\
R^0&=\frac{u}{r'w}\\
S^0&=\frac{t't'}{s'}.\frac{v}{w}
\end{aligned}
\]
The fourteen formulas I, II, III give the complete solution of our problem in the simplest possible form.

It is also worth noting that for $r'=r''$, whether the comparison observations give this equality or whether, in their absence, we are unable to take into account the variability of $R$ during the two sets of observations, we have
\[
\begin{aligned}
V&=\tfrac12(B'-B''-q'+q'')\pm90^0\\
v&=\pm2\sin\tfrac12(B'-B''-q'+q'')
\end{aligned}
\]
where the upper or lower signs apply according as
\[
\sin\tfrac12(B'-B''-q'+q'')
\]
is positive or negative.

\Article{8}
To illustrate these rules we add the complete calculation of an example. The calculation is carried out with seven-figure logarithms. This is much more precise than is needed in practice, where five-figure logarithms always suffice.

On 24 March 1841 the swing duration of the bifilar magnetometer was found from observations extending over $1^{St}21'$ (after the appropriate reduction to infinitesimally small swings, of course) to be $=28''89071=t'$. The setting of the mirror alidade was $154^020'30''=A'$, and that of the cradle alidade was $=27^040'25''=B'$. The mean position from several determinations evenly distributed over that interval was $994.33$ scale divisions. Since the plumb line corresponds to the scale reading $1000$ and a scale division amounts to $21''5835$, we have $p'=-2'2''38$. From exactly simultaneous observations, the swing duration of the unifilar magnetometer at the magnetic
\SourceLeaf{425}observatory was found to be $=20''72725$ and the mean position $=881.80$ scale divisions. The mean of the daily records for February, $888.40$, was chosen as the reference position (corresponding to the absolute declination $18^011'54''$). Since a scale division on the unifilar magnetometer amounts to $21''3489$, we obtain $q'=-2'20''90$.

The mean thermometer readings (from records taken immediately before the observations began and immediately after they ended) were $+6^096$ in the case of the bifilar magnetometer, $+7^06$ at the upper attachment of the suspension wires, and $+7^045$ in the case of the unifilar magnetometer, all on the Réaumur scale.

Likewise, for a second set of observations on the following day we had
\[
\begin{aligned}
t''&=108''17\\
A''&=151^027'30''\\
B''&=185\;59\;35\\
p''&=-24'33''07\\
q''&=+2\;42.04
\end{aligned}
\]
the swing duration of the unifilar magnetometer $=20''73117$, and the thermometer readings in the same order as above $+6^036$, $+7^00$, $+7^01$.

To calculate the effect of temperature, I set
\[
\alpha=0.000016,\quad\beta=0.000024,\quad\gamma=\gamma'=0.000765
\]
taking the last value from \textsc{Hansteen}, since I still lack conclusive determinations of my own. From formulas (1) and (3) of Article 5 it follows, if we choose $20''72=\theta^0$ as the reference swing duration of the unifilar magnetometer, that
\[
\begin{aligned}
\log r'&=-0.0001376\\
\log s'&=-0.0000043\\
\log r''&=-0.0002155\\
\log s''&=-0.0000044
\end{aligned}
\]
Using the abbreviated formulas (2) and (4) of the same article, we could set
\[
\begin{aligned}
\log s&=-\alpha k(c''-c)\\
\log r&=2\log\frac{\theta^0}{\theta}+(\gamma+2\alpha)k(c'-c)-\alpha k(c''-c)
\end{aligned}
\]
\SourceLeaf{426}if $k$ denotes the modulus of the Briggs logarithms. Thus, with the values of $\alpha,\beta,\gamma$ given above,
\[
\begin{aligned}
\log s&=0.00000695(c''-c)\\
\log r&=2\log\frac{\theta^0}{\theta}+0.0003461(c'-c)-0.00000695(c''-c)
\end{aligned}
\]
from which we obtain for our observations
\[
\begin{aligned}
\log r'&=-0.0001386\\
\log s'&=-0.0000044\\
\log r''&=-0.0002169\\
\log s''&=-0.0000044
\end{aligned}
\]
values whose difference from those above is barely perceptible.

After these preparations, the main steps of the calculation itself are as follows
\[
\begin{aligned}
A'-A''-p'+p''&=+2^030'29''31\\
B'-B''-q'+q''&=-158\;14\;7.06\\
\log\frac{t't's''}{s't''t''}&=8.8524525\\
\log\frac{r''}{r'}&=9.9999221
\end{aligned}
\]
From these, by I,
\[
\begin{aligned}
\log\mathfrak{A}&=8.8520363\\
\log\mathfrak{A}_1&=7.4935432\\
\log\mathfrak{B}&=9.9678043\,\mathrm{n}\\
\log\mathfrak{B}_1&=9.5690569\,\mathrm{n}
\end{aligned}
\]
from which further
\[
\begin{aligned}
\log(\mathfrak{A}-1)&=9.9679562\,\mathrm{n}\\
\log(\mathfrak{B}-1)&=0.2852304\,\mathrm{n}\\
\log(\mathfrak{B}-\mathfrak{A})&=9.9998589\,\mathrm{n}\\
\log(\mathfrak{B}_1-\mathfrak{A}_1)&=9.5726915\,\mathrm{n}
\end{aligned}
\]
Then formulas II give
\SourceLeaf{427}
\[
\begin{aligned}
U&=179^048'28''15\\
V&=190\;52\;52.91\\
W&=200\;30\;14.79\\
\log u&=9.9679586\\
\log v&=0.2931101\\
\log w&=0.0282829
\end{aligned}
\]
and finally formulas III give
\[
\begin{aligned}
P^0&=144^045'10''50\\
Q^0&=7\;0\;59.26\\
\log R^0&=9.9398133\\
\log S^0&=3.1863476
\end{aligned}
\]

\Article{9}
The problem can be generalized further by taking \textit{four different} sets of observations as the basis, two for the position and two for the swing duration, while dropping the assumption that the sets of either kind correspond respectively to the same values of $B$. This has the advantage that the position of the bifilar magnetometer can be observed with the bar magnet almost entirely at rest, using the procedure described in the \textit{Resultate} for 1836. II. [see below]. That advantage loses its value, however, because we still have to observe both the swing duration and the position on the auxiliary magnetometer for all four sets, and therefore still have to determine the latter from the elongations. It is clear, then, that this method takes twice as much work as the one in Article 7, which also has the advantage of a very simple calculation. Directly determining the quantities from four separate sets of observations is much more involved, and we would therefore prefer to reserve its development, which is of some mathematical interest, for another place.

\Article{10}
It is worth noting what happens when one determines the position from observed elongations by applying the procedure [see below] in the \textit{Resultate} 1836 II
\SourceLeaf{428}without modification. The unequal angular values of the scale divisions cause an error. The farther the position is from the centre of the scale, the greater the error. If fully precise results are required, the procedure must therefore be applied to the distances of the elongations from the centre of the scale after converting them into divisions of arc by an exact formula, rather than directly to the scale divisions read at the elongations. If the position is near the centre, the error is of course immaterial. We shall therefore always choose the setting of the mirror, or the value of $A$, so that the position differs little from the centre, or so that $p$ is small. When the quantities are first determined, this can of course be achieved only by a preliminary trial (in much the same way as in Article 6). If the quantities $P,Q,R$ are already known approximately, however, a calculation is preferable for this purpose. It is most conveniently given the following form, which is easily derived from Article 2, formula (1) or (4). Determine an angle $\psi$ from the formula
\[
\operatorname{tang}\psi=\frac{1-R}{1+R}.\operatorname{tang}\tfrac12(Q-B)=\operatorname{tang}\tfrac12\varphi^2.\operatorname{tang}\tfrac12(Q-B)
\]
choosing $\psi$ in the same quadrant as $\tfrac12(Q-B)$, and then set
\[
A=\psi+P-\tfrac12(Q-B)
\]

\Article{11}
It remains to establish how the observations with the bifilar magnetometer in the transverse position are related to changes in the quantities.

We assume that the transverse position determined by the procedure of Article 3 refers to the reference values of the quantities. The cradle is thus set so that, at rest, the magnetic axis of the bar magnet makes a right angle with the reference magnetic meridian when the ratio of the magnetic directive force to the static directive force is $R^0$ to $1$. The mirror is set so that in this position the image of the point on the scale marked by the plumb line appears at the crosshairs of the observing telescope. Thus, if we denote the swing duration under these conditions by $t^0$, we have
\SourceLeaf{429}
\[
\begin{aligned}
\sin\varphi&=R^0\\
A&=P\mp\varphi\\
B&=Q^0\mp(90^0+\varphi)\\
\frac{S^0}{t^0t^0}&=\cos\varphi
\end{aligned}
\]
where the paired signs refer to the westward or eastward position of the north pole of the bar magnet. If we now retain the symbols
\[
p,\quad Q=Q^0+q,\quad R=rR^0,\quad S=sS^0,\quad t
\]
in their previous general sense, formulas (1) and (6) of Article 2 give
\[
\begin{aligned}
\sin(\varphi\pm p)&=rR^0\cos(p-q)\\
\frac{sS^0}{tt}\sin(\varphi\pm p)&=rR^0\cos(\varphi\pm q)
\end{aligned}
\]
or
\[
\begin{aligned}
r&=\frac{\sin(\varphi\pm p)}{\sin\varphi\cos(p-q)}\quad .\ .\ .\ .\ .\ .\ .\ .\ .\quad(1)\\
\frac{st^0t^0}{tt}&=\frac{\cos(\varphi\pm q)}{\cos\varphi\cos(p-q)}\quad .\ .\ .\ .\ .\ .\ .\ .\ .\quad(2)
\end{aligned}
\]

\Article{12}
The most important use of the bifilar magnetometer is to determine changes in horizontal intensity. These are related to changes in $R$ by the formula given above (Article 5)
\[
r=\frac{f^0g^0hMT}{fgh^0M^0T^0}
\]
Recall that $T^0$ denotes the reference intensity chosen at the outset and $M^0$ denotes the magnetic moment of the bar magnet reduced to the reference temperature, according to the magnetic state of the bar when the constants were determined. If we denote the magnetic moment at an unspecified time, likewise reduced to the reference temperature, by $\mathfrak{M}$ and set
\[
\frac{\mathfrak{M}}{M^0}=\lambda,\qquad\frac{T^0}{\lambda}=\mathfrak{T}
\]
then, under the conditions discussed in Article 4, $\lambda$ is a coefficient that can be regarded as constant for a moderate period of
\SourceLeaf{430}time, say one day. As it and $\mathfrak{M}$ gradually decrease very slowly, $\mathfrak{T}$ will gradually increase and will always denote the horizontal intensity at which, at the reference temperature, the ratio of the magnetic directive force to the static directive force equals the ratio $R^0:1$. The equation above now takes the form
\[
T=\frac{\mathfrak{M}fgh^0}{Mf^0g^0h}.r\mathfrak{T}
\]
where the first factor $\frac{\mathfrak{M}fgh^0}{Mf^0g^0h}$ depends only on the temperature and (if we retain the notation of Article 5) can be expressed by
\[
1+(\gamma+2\beta-\alpha)c+(\theta-\tfrac12\alpha)(c''-c)
\]
Meanwhile $r$ can be determined at any instant from formula (1) of the preceding article by combining simultaneous position observations with the bifilar magnetometer in the transverse position for $p$ and with the unifilar magnetometer for $q$. It is therefore clear that changes in intensity can be followed with the greatest precision over the shortest intervals in this way, so long as we only need to compare the changed intensities \textit{with one another} over a moderate period, for example during a twenty-four-hour scheduled observation session or during the time required for an absolute determination of intensity with the unifilar magnetometer. If, in such an absolute determination, we use simultaneous observations with the bifilar magnetometer to reduce the individual operations to a single reference intensity (see \textit{Intensitas vis magneticae}, Articles 10 and 22), we also obtain the value of $\mathfrak{T}$ valid for that time in absolute measure. This is in itself preferable to using a second unifilar magnetometer as recommended there. If we repeat such absolute determinations from time to time, we can keep track of any gradual changes in $\mathfrak{T}$ and estimate those changes for the intervening time by interpolation in geometric progression without appreciable error. In this way we can give all changes in intensity, in all their variations, in absolute measure. It goes without saying that if, after a longer interval, secular changes in magnetic declination and horizontal in%
\SourceLeaf{431}tensity or a considerable weakening of the bar’s magnetism cause $p$ and $q$ to cease to stay within moderate limits (cases of large, exceptional anomalies need not count here), we shall make a suitable adjustment to the setting of the cradle and the mirror and, if we consider it advisable, to the separation of the suspension wires. We shall then begin a new series of observations with changed quantities.

\Article{13}
As long as $p$ and $q$ are small, for any purpose that does not require the greatest precision we can replace the exact formula (1) with a shorter one in which $q$ drops out entirely. There is then no need for simultaneous observations with the unifilar magnetometer. This applies in particular to ordinary scheduled observations. In place of that formula we can set
\[
r=1\pm\operatorname{cotang}\varphi.\operatorname{tang}p\quad\text{or also}\quad r=1\pm\tfrac12\operatorname{cotang}\varphi.\operatorname{tang}2p
\]
Now, if $n$ denotes the difference between the scale division read and the one to which the plumb line refers, and $d$ denotes the horizontal distance from the centre of the mirror to the latter point, measured in scale divisions, then
\[
\operatorname{tang}2p=\frac{n}{d}
\]
so the formula becomes
\[
r=1\pm\frac{n}{2\operatorname{tang}\varphi.d}
\]
and at the same time, with sufficient accuracy,
\[
T=\mathfrak{T}(1\pm\frac{n}{2\operatorname{tang}\varphi.d}+(\gamma+2\beta-\alpha)c)
\]
if we omit the small term $(\beta-\tfrac12\alpha)(c''-c)$. With the apparatus used here, $d=4778{,}3$ millimetres, and the results of the calculation given as an example in Article 8 yield $\varphi=60^031'37''9$, hence $2d\operatorname{tang}\varphi=16910$. With the values of $\alpha,\beta,\gamma$ used there we thus obtain
\[
T=(1+\frac{n+13{,}65c}{16910})\mathfrak{T}
\]
\SourceLeaf{432}if the north pole of the bar magnet is on the west side, and
\[
T=(1-\frac{n-13{,}65c}{16910})\mathfrak{T}
\]
if it is on the east side.

It goes without saying that temperature can safely be disregarded altogether in scheduled observations so long as the only aim is to discern the shape of the individual anomalies that change over short intervals.

\Article{14}
Just as in the transverse position of the bifilar magnetometer changes in intensity affect the position with their full strength and changes in declination barely affect it, so the reverse holds for the swing duration. The latter changes have the greatest effect on it, and the former have only an extremely small effect. Insofar as $p,q$ and the deviation of the quantity $S$ from its reference value are small, we can, without appreciable error, replace formula (2) of Article 11 with
\[
\frac{t^0t^0}{tt}=1\mp q\operatorname{tang}\varphi\quad\text{or also}\quad t=t^0(1\pm\tfrac12q\operatorname{tang}\varphi)
\]
when $q$ is expressed in units of the radius, and consequently with
\[
t=t^0\pm\frac{t^0\operatorname{tang}\varphi}{412530}.q
\]
when it is expressed in seconds of arc. The results of the example calculated above give $t^0=55{,}871$ seconds of time, so in seconds of arc we have
\[
q=\pm(t-55{,}871)4172''8
\]
The exact transformation of formula (2) for calculating $q$ is as follows
\[
\pm\operatorname{tang}q=\frac{tt-st^0t^0\cos p}{tt\operatorname{tang}\varphi\pm st^0t^0\sin p}
\]
It goes without saying that swing obser\-vations cannot determine the value of $q$ at a particular instant this way and can determine only its mean value over the duration of those observations.

